% Generated draft: requires mathematical and source review.
\documentclass[11pt]{article}
\usepackage[margin=25mm]{geometry}
\usepackage{fontspec}
\setmainfont{Latin Modern Roman}
\usepackage{amsmath,amssymb,mathtools,unicode-math}
\setmathfont{Latin Modern Math}
\usepackage{enumitem,needspace,xurl,hyperref}
\setlength{\parindent}{0pt}
\setlength{\parskip}{5pt}
\setlength{\emergencystretch}{3em}
\newcommand{\N}{\mathbb{N}}
\newcommand{\Z}{\mathbb{Z}}
\newcommand{\Q}{\mathbb{Q}}
\newcommand{\R}{\mathbb{R}}
\newcommand{\C}{\mathbb{C}}
\newcommand{\F}{\mathbb{F}}
\newcommand{\A}{\mathbb{A}}
\newcommand{\PP}{\mathbb{P}}
\newcommand{\E}{\mathbb{E}}
\newcommand{\eps}{\varepsilon}
\newcommand{\dd}{\,\mathrm{d}}
\begin{document}
\section*{UnsolvedMath ID: 11100020}
\textbf{Source problem code:} catalogue entry
\subsection{Definitions and mathematical statement}
Fix a prime $p$, which the source does not specify. Write $K(1)$ for first Morava $K$-theory at $p$, $E(1)$ for the height-one Johnson--Wilson theory, and $L_{K(1)}$ and $L_{E(1)}$ for Bousfield localization of spectra with respect to these theories. A spectrum $X$ is $K(1)$-local if the localization unit $X\to L_{K(1)}X$ is an equivalence. Let $\operatorname{Ho}(\mathrm{Sp}_{K(1)})$ be the category of such spectra, with stable homotopy classes of maps as morphisms.

The source asks for a more natural description of Bousfield's algebraic description of the $E(1)$-local category when considered in the $K(1)$-local setting. In particular, seek an explicit algebraic description of $\operatorname{Ho}(\mathrm{Sp}_{K(1)})$, explaining the invariants of its objects and the information needed to recover morphisms and their composition. The source motivates this request by reporting that Bousfield describes isomorphism classes but not maps. This motivation does not itself specify a particular algebraic target category, equivalence, or required notion of naturality; these remain part of the formulation to be clarified. The source gives no restriction to finite or dualizable spectra and does not fix an odd-prime hypothesis.

Bousfield's description and Franke's claimed generalization are historical context in the original item. The supplied later summary reports substantial height-one progress, but its assertion that equivariant homotopy modules classify both objects and maps is not supplied with a precise theorem or hypotheses. It is therefore not used here as a categorical classification theorem.
\subsection{Short English statement}
Find a natural algebraic account of the $K(1)$-local category that explains Bousfield's work and clarifies the treatment of maps as well as objects. The original author reports difficulty understanding Franke's claimed generalization and says that no improvement on Bousfield had succeeded in the preceding ``10 plus years''; this is a source-era assessment. The registry labels the record partially solved and reports that height one is largely understood while a higher-height program remains open. These reports are not a verified current classification result, and the exact meaning of ``more natural'' remains unspecified.
\subsection{Sources}
\begin{itemize}
\item Local UnsolvedMath registry, permanent ID $11100020$, source code \texttt{catalogue entry}, proposed by Mark Hovey; record created 2026-07-29 and updated 2026-08-06. Supplied status: \texttt{partially\_solved}; research classification: \texttt{PARTIAL-PROGRESS}.
\item Mark Hovey, \emph{Hovey's Algebraic Topology Problem List}, ``Morava K- and E-theory,'' item 10. Surviving mirror: \url{https://www-users.cse.umn.edu/~tlawson/hovey/}, accessed 2026-07-29; extraction method: HTML ordered list. The background describes 65 top-level items in nine category pages, assigns default difficulty L3, and marks status evidence \texttt{NEEDS\_REVIEW}: the source presents an open problem and a current primary-literature status audit remains outstanding.
\item The statement names Bousfield's $E(1)$-local description and Franke's claimed generalization, without bibliographic particulars. The research summary additionally invokes Morava/Dieudonn\'e theory, Bousfield, Strickland, and Barthel--Heard--Sanders, mentioning $K(1)$ and an odd-primary ``PBSS'' without defining that acronym or giving exact references. It describes odd-primary objects and maps through $\Z_p^{\times}$-equivariant $\pi_*$-modules, mentions a finiteness condition and a spanning assertion without precise scope, and describes the prime-two case as subtler because of $2$-adic units and $\eta,\nu$. Its final ``Franke's claim'' discussion is truncated by an ellipsis. All these are supplied summary claims requiring verification.
\item Dataset attribution: Ulam AI / UnsolvedMath Contributors, CC BY 4.0. Separate source-rights note: \texttt{NEEDS\_REVIEW}; publicly accessible author/publisher copy, with redistribution terms requiring review.
\end{itemize}
\end{document}
